\chapter{State of the Art and Taxonomy}
\label{ch:master-sota}

\section{Introduction}

This chapter applies the structured mapping protocol defined in
\cref{ch:master-background} and \cref{app:master-review-protocol} to the
Android third-party library detection literature. Its purpose is not to
produce a numerical leaderboard. The reviewed studies use different datasets,
task units, obfuscation settings, and ground-truth procedures, so a direct
ranking would hide the very uncertainty that this thesis aims to expose.
Instead, the chapter organizes prior work by technique, evidence source,
transformation threat, and evaluation practice.

The analysis is centered on the 15 Tier A studies whose full texts are
available locally. Tier B records are useful for search coverage and for
identifying emerging directions, especially SBOM-oriented recovery and recent
benchmarking, but they are not used here as independent authority for precise
quantitative claims. This boundary keeps the review consistent with the
protocol and prevents unverified metadata from receiving the same weight as
inspected studies. Fields not reported by a Tier A publication are treated as
not reported rather than inferred, and original-paper claims are kept
separate from later cross-tool observations.

The chapter uses the term \emph{component recovery} broadly. A study may
recover a library family, a package cluster, a class set, an exact version, a
vulnerability-related version list, or only a presence signal. For trustworthy
\SBOM reconstruction, these outputs are not equivalent. A method that detects
the presence of a library family does not necessarily justify an exact Maven
version, and an exact-version candidate does not automatically justify a
\CVE claim unless the evidence, reference coverage, and validation context are
visible.

\section{Taxonomy of Component-Recovery Techniques}

The reviewed approaches can be grouped into seven recurring technique
families. The categories are not mutually exclusive: most mature tools combine
several families, for example by using a fast coarse filter followed by a
slower structural match. The taxonomy therefore describes the role each
family plays in a recovered component claim.

\begingroup
\small
\renewcommand{\arraystretch}{1.15}
\begin{longtable}{p{0.20\textwidth}p{0.22\textwidth}p{0.28\textwidth}p{0.18\textwidth}}
\caption{Technique taxonomy for Android component recovery.}
\label{tab:master-technique-taxonomy}\\
\toprule
\textbf{Technique family} & \textbf{Typical evidence} & \textbf{Trust role} & \textbf{Tier A examples}\\
\midrule
\endfirsthead
\toprule
\textbf{Technique family} & \textbf{Typical evidence} & \textbf{Trust role} & \textbf{Tier A examples}\\
\midrule
\endhead
\bottomrule
\endfoot
Namespace and metadata heuristics &
Package trees, class names, library naming conventions &
Cheap candidate generation, but weak authority when names are renamed,
flattened, relocated, or shaded. &
Assessed in \cite{zhan2020assessment,zhan2021systematic}\\

Profile, hash, and signature matching &
Method descriptors, Merkle trees, fuzzy hashes, class or method signatures &
Can support identity and, when references are complete, version claims. Its
authority is bounded by reference coverage and by the stability of the hashed
features under optimization. &
LibScout \cite{backes2016libscout}, ORLIS \cite{wang2018orlis}, LibID
\cite{zhang2019libid}\\

Structural graph segmentation &
Call, inheritance, containment, and class-dependency relations &
Improves module boundary recovery when package structure is unreliable.
Supports stronger identity evidence, but does not by itself solve exact
version attribution. &
LibD \cite{li2017libd}, LibPecker \cite{zhang2018libpecker}, assessment
studies \cite{zhan2020assessment,zhan2021systematic}\\

Coarse-to-fine matching &
Fast filtering followed by detailed opcode, CFG, or structural comparison &
Balances scalability and precision. Trust depends on both stages: a missed
coarse candidate can cap recall, while an over-permissive fine stage can
inflate false positives. &
ORLIS \cite{wang2018orlis}, ATVHUNTER \cite{zhan2021atvhunter}, LibScan
\cite{wu2023libscan}\\

Version-discriminating semantic features &
Opcode sequences, basic-block features, call chains, field-operation
patterns, literals, and class summaries &
Provides the main route from library identity to exact-version authority.
Adjacent versions and optimized builds remain difficult because many versions
share most code. &
ATVHUNTER \cite{zhan2021atvhunter}, LibScan \cite{wu2023libscan}, LibHunter
\cite{xie2024libhunter}, SAD \cite{zhou2025sad}\\

Learned representations &
Smali token sequences, structural vocabularies, neural embeddings &
Can learn obfuscation-resilient presence signals, but the reviewed Tier A
learning approach does not provide exact-version or boundary authority on its
own. &
LibAttention \cite{pan2025libattention}\\

Empirical assessment and benchmarking &
Shared datasets, patched tools, source-known applications, evaluation
criteria &
Does not directly recover components, but defines which claims are comparable
and exposes tool failure modes, missing denominators, and reproducibility
limits. &
Tool assessments \cite{zhan2020assessment,zhan2021systematic}, Ramirez
thesis \cite{ramirez2024empirical}\\
\end{longtable}
\endgroup

The movement across these families is historically important. Early methods
made library profiles and package structure usable at scale, but later
studies increasingly treat package names as unreliable evidence. Graph-based
and version-aware methods improve the situation by using relations, opcode
content, and class behavior. The newest verified Tier A work adds learned
representations and finer class-level summaries. Even so, the taxonomy shows
that standard SBOM emission remains outside the main output model of the
verified Android TPL literature: most studies output names, clusters,
versions, class mappings, or vulnerability-linked lists rather than a
provenance-rich CycloneDX or SPDX document.

\section{Taxonomy of Evidence Sources}

An evidence source is useful for trustworthy \SBOM recovery only when its
origin, stability, version specificity, and provenance are understood. The
same signal may be strong for identity and weak for exact version. Conversely,
a highly version-specific literal may be too fragile to survive resource
removal or string encryption. \Cref{tab:master-evidence-taxonomy} summarizes
the recurring evidence families extracted from the Tier A studies.

\begingroup
\small
\renewcommand{\arraystretch}{1.15}
\begin{longtable}{p{0.18\textwidth}p{0.23\textwidth}p{0.21\textwidth}p{0.28\textwidth}}
\caption{Evidence-source taxonomy for reconstructed Android component claims.}
\label{tab:master-evidence-taxonomy}\\
\toprule
\textbf{Evidence family} & \textbf{Examples} & \textbf{Typical stability} & \textbf{Trust implication}\\
\midrule
\endfirsthead
\toprule
\textbf{Evidence family} & \textbf{Examples} & \textbf{Typical stability} & \textbf{Trust implication}\\
\midrule
\endhead
\bottomrule
\endfoot
Names and namespaces &
Package prefixes, class names, manifest names &
Weak under renaming, package flattening, relocation, and shading &
Useful for fast hints and human inspection, but insufficient for exact SBOM
claims without corroboration.\\

Descriptors and API surfaces &
Method signatures, fuzzy type descriptors, Android API usage &
More stable than names under identifier renaming; weaker under dead-code
removal and inlining &
Supports identity and coarse version discrimination when combined with
reference profiles.\\

Opcode and instruction sequences &
Basic-block opcodes, method opcode chains, call-chain opcodes &
Resilient to naming changes; threatened by optimization, inlining, and
control-flow rewriting &
Provides version-discriminating evidence, especially when compared against
known library releases.\\

Control-flow and dependency structures &
CFGs, call graphs, inheritance graphs, class-dependency graphs &
Stable against simple renaming and often against package restructuring;
affected by optimization and partial inclusion &
Improves module boundary recovery and reduces dependence on package layout.\\

Literals, constants, and resources &
Strings, constants, resource identifiers, assets &
Potentially version-specific, but fragile under stripping, string encryption,
and resource rewriting &
Best used as tie-breakers or independent corroboration, not as sole authority.\\

Reference corpus and release metadata &
Library profiles, Maven coordinates, release chronology, known vulnerable
version ranges &
External to the \APK; stable only if the corpus is complete and correctly
constructed &
Defines the ceiling for recognizable identities and versions. Missing
references must not be interpreted as component absence.\\

Learned representations &
Transformer embeddings over Smali tokens, structural vocabularies &
Can absorb complex obfuscation patterns, but may be hard to audit feature by
feature &
Promising for presence screening and robustness, but provenance and exact
version authority require additional explanation.\\

Validation evidence &
Controlled builds, source manifests, F-Droid dependency data, source-known
ground truth &
Strong within the validation context; not automatically generalizable &
Can upgrade a claim's evaluation status, but must remain separate from
APK-only inference.\\
\end{longtable}
\endgroup

The trust implication of this taxonomy is that evidence fusion should prefer
independent and complementary evidence over repeated variants of the same
feature. For example, several hash matches derived from the same method body
are less independent than a method-level match corroborated by a
class-dependency relation and a version-specific literal. This distinction
supports the thesis synthesis: independence, provenance, and reference
coverage must be treated as first-class quality dimensions whenever detector
outputs are interpreted as SBOM claims.

\section{Taxonomy of Transformation Threats}

Android component recovery is difficult because compilation and packaging can
remove the very boundaries that build-time dependency tools rely on. R8 and
related obfuscators may rename identifiers, flatten packages, remove unused
code, rewrite control flow, and inline methods \cite{androidShrinkCode}.
The reviewed literature shows that no single evidence family remains stable
under all transformations.

\begingroup
\small
\renewcommand{\arraystretch}{1.15}
\begin{longtable}{p{0.18\textwidth}p{0.25\textwidth}p{0.22\textwidth}p{0.25\textwidth}}
\caption{Transformation-threat taxonomy for Android component recovery.}
\label{tab:master-transformation-taxonomy}\\
\toprule
\textbf{Threat} & \textbf{Evidence most affected} & \textbf{More resilient evidence} & \textbf{Trust consequence}\\
\midrule
\endfirsthead
\toprule
\textbf{Threat} & \textbf{Evidence most affected} & \textbf{More resilient evidence} & \textbf{Trust consequence}\\
\midrule
\endhead
\bottomrule
\endfoot
Identifier renaming &
Class, method, field, and package names &
Fuzzy descriptors, opcode features, structural relations &
Name matches cannot be treated as independent authority unless corroborated.\\

Package flattening and relocation &
Package-subtree segmentation and namespace heuristics &
Class-dependency graphs, call relations, profile matching independent of
package roots &
Package-based module boundaries become unreliable; graph segmentation becomes
more important.\\

Control-flow rewriting &
CFG shape, branch structure, fuzzy CFG hashes &
Higher-level dependency relations, method descriptors, class summaries &
CFG-only methods may overstate robustness when control-flow transformations
are not in the evaluation.\\

Dead-code removal and shrinking &
Presence of optional classes, methods, API calls, literals &
Partial-match scoring, class-level confidence, source-known validation &
Absence of a feature may indicate transformation rather than absence of the
library.\\

Method inlining and outlining &
Method boundaries, call-chain evidence, version-specific method bodies &
Multi-scale evidence, class-level summaries, literals, residual dependency
structure &
Exact-version authority must be capped when inlining changes the features
used for version discrimination.\\

Multidex and split packaging &
Single-file DEX assumptions and incomplete scanners &
Whole-APK extraction over all DEX files and split artifacts &
Evaluation denominators must state whether all executable artifacts were
analyzed.\\

Partial library inclusion &
Binary present/absent decisions and complete-profile matching &
Similarity scores, per-class evidence, explicit ambiguity classes &
A partial match should not automatically become a full component/version
claim.\\

Nested or transitive dependencies &
Flat library lists and one-library-one-package assumptions &
Dependency graphs, source-assisted validation, relationship-aware SBOM
models &
A recovered component may be direct, transitive, bundled, shaded, or copied;
the relationship must remain uncertain unless evidenced.\\

Native and dynamic code &
DEX-only static features &
Native file hashes, load-site evidence, dynamic analysis, source validation &
Out-of-scope code must be represented as a boundary, not silently ignored.\\
\end{longtable}
\endgroup

The most consequential threats for exact-version claims are shrinking,
partial inclusion, and method inlining. Identifier renaming is comparatively
well handled by modern structural or opcode-based approaches, but version
evidence often depends on small differences between adjacent releases.
Optimizations that delete, merge, or inline exactly those differences can turn
an apparently confident match into a version-band ambiguity.

\section{Representative Android TPL Approaches}

\subsection{From Profile Matching to Structural Boundaries}

LibScout established a reference-profile model for matching compiled Android
applications against known library releases. Its three-layer profile structure
and fuzzy method descriptors made identifier renaming less damaging, and its
similarity scores provided a graded alternative to exact profile equality
\cite{backes2016libscout}. From the viewpoint of trustworthy \SBOM recovery,
the important lesson is dual: reference profiles can support version claims,
but those claims inherit the assumptions of the profile structure and the
reference corpus. When package layout or method boundaries are transformed,
the confidence of the resulting claim must be reduced rather than silently
preserved.

LibD moved away from pure package-layout dependence by recovering library
instances through a homogeny graph over inclusion, inheritance, and call
relations, combined with opcode-derived fingerprints \cite{li2017libd}.
This made large-scale library discovery more resilient to naming changes, but
the output remained a clustering-oriented library identification result
rather than an exact-version or SBOM result. For the taxonomy, LibD is
therefore strong evidence for graph-based boundary recovery and weak evidence
for version authority.

ORLIS introduced a two-stage obfuscation-resilient design. Its first stage
uses interprocedural call-graph-derived features and \texttt{sdhash} to remove
unlikely library candidates. Its second stage applies \texttt{ssdeep} class
signatures and an injective class mapping to verify whether a library is
included in the application \cite{wang2018orlis}. The design is valuable
because it separates scalable filtering from detailed confirmation and can
represent partial inclusion without relying on package names. Its limitations
are equally important for the proposed framework: a class mapping is not
automatically an exact-version claim, and computational cost affects whether a
method can be used as routine SBOM evidence.

LibPecker and LibID further illustrate the shift from names to structural and
class-level evidence. LibPecker emphasizes class dependency and adaptive
similarity. LibID formulates consistent class correspondence with binary
integer programming and identifies specific library versions under code
shrinking and package modification
\cite{zhang2018libpecker,zhang2019libid}. Together with the later assessment
studies, they support the review's central classification: package structure
is a fragile feature family, while class relationships, signatures, and
opcode-derived evidence are more suitable as identity support under
obfuscation.

\subsection{Version-Aware and Vulnerability-Oriented Pipelines}

ATVHUNTER combines class-dependency based decoupling, coarse CFG filtering,
fine-grained opcode matching, and a vulnerability database to identify
vulnerable library versions \cite{zhan2021atvhunter}. It is one of the most
important verified studies for this thesis because it connects version
identification to security interpretation. However, its output remains a
version and vulnerability list rather than a standard \SBOM with explicit
provenance, uncertainty, and claim boundaries. The study therefore motivates
the framework's rule that vulnerability authority cannot exceed the strength
of the underlying version evidence.

LibScan and LibHunter extend the version-aware family by using richer opcode,
class-correspondence, call-chain, and optimization-aware evidence
\cite{wu2023libscan,xie2024libhunter}. Their relevance is not only that they
improve matching. They also show that code optimization itself is an
evaluation dimension. In particular, LibHunter augments reference features to
model call-site optimization and performs cross-inlining method matching. A
method evaluated only on clear or lightly obfuscated applications has not
demonstrated the same authority as a method evaluated under inlining and
shrinking.

SAD provides the most detailed Tier A example of class-structural
version-level reasoning. It splits classes into stateful and stateless
categories, applies field-operation sequences or opcode similarity
accordingly, and uses literals to break ties between adjacent versions
\cite{zhou2025sad}. This design supports the thesis claim that exact-version
attribution requires evidence at a finer granularity than library presence.
At the same time, its remaining sensitivity to method inlining confirms that
strong reported version-level performance does not eliminate transformation
risk.

\subsection{Learning-Based Presence Detection}

LibAttention represents a different family: instead of comparing explicit
profiles, it learns representations from large Smali corpora and predicts
library presence with a Transformer-based model
\cite{pan2025libattention}. This is important because learned features can
capture patterns that rule-based signatures miss under aggressive
obfuscation. For \SBOM reconstruction, however, the reviewed Tier A result is
bounded: it predicts presence rather than exact version, does not expose an
explicit class or library-boundary mapping, and does not provide standard SBOM
output. In the taxonomy, learned presence detection is therefore a potentially
useful candidate-generation signal but not sufficient version or vulnerability
authority by itself.

\subsection{Empirical Assessments and Cross-Tool Studies}

The ASE 2020 assessment and the later systematic assessment are not merely
background surveys; they define several trust requirements for any recovered
\SBOM pipeline \cite{zhan2020assessment,zhan2021systematic}. They show that
tool results depend on module decoupling, multidex handling, AAR processing,
obfuscation setting, version-identification capability, efficiency, and ease
of reproduction. These studies are a warning against reporting a single
accuracy value without denominators, exclusions, and task levels.

The Ramirez thesis provides a cross-platform empirical comparison of Android
and iOS library-detection tools \cite{ramirez2024empirical}. Its Android study
reuses 75 APKs from the ORLIS dataset, 75 from the LIBLOOM dataset, and 50
manually selected Google Play applications. It is therefore useful for
cross-tool execution and dataset-reuse lessons, but it is not a source-known
or F-Droid validation study and it does not independently establish standard
SBOM output.

TPL-Benchmark strengthens the benchmark perspective with a large,
version-annotated Android dataset and a standardized reevaluation of ten
detection methods \cite{gu2025tplbenchmark}. Its compatible tool-evaluation
slice shows that the best library-level F1 score was 60.15\%, while LibHunter
obtained the highest reported version-level recall at 76.48\%. These results
are not a universal ceiling, but they demonstrate that realistic shared
ground truth can produce substantially weaker and more heterogeneous outcomes
than isolated original evaluations.

UniBOM is an adjacent rather than Android-specific study. It combines firmware
extraction, filesystem scanning, source-code analysis, SBOM generation, and
vulnerability visualization, and evaluates the workflow on 258 wireless
gateway firmware images and four IoT operating-system codebases
\cite{safronov2025unibom}. It confirms that binary, metadata, and source
evidence can coexist in one SBOM workflow. It does not, however, solve
APK-side Maven exact-version attribution or provide ground truth for Android
TPL detection.

\section{Evaluation Practices and Ground Truth}

The reviewed studies evaluate different questions even when they use similar
words such as detection, identification, precision, recall, or version
accuracy. A trustworthy synthesis must therefore classify evaluation practice
before interpreting the metric.

\begingroup
\small
\renewcommand{\arraystretch}{1.15}
\begin{longtable}{p{0.22\textwidth}p{0.27\textwidth}p{0.39\textwidth}}
\caption{Evaluation-practice taxonomy in the Android TPL literature.}
\label{tab:master-evaluation-taxonomy}\\
\toprule
\textbf{Evaluation practice} & \textbf{What it can show} & \textbf{Boundary for SBOM trust}\\
\midrule
\endfirsthead
\toprule
\textbf{Evaluation practice} & \textbf{What it can show} & \textbf{Boundary for SBOM trust}\\
\midrule
\endhead
\bottomrule
\endfoot
Controlled library or reference corpus &
Whether a matcher can distinguish known library profiles or versions under
specified transformations &
Results depend on reference completeness and may overstate performance on
diverse real applications.\\

Synthetic or tool-based obfuscation tests &
Whether selected transformations damage a feature family &
Only covers the transformations and configurations tested; R8 inlining and
shrinking require explicit treatment.\\

Market-scale discovery &
Whether a detector scales and finds plausible library families in many apps &
Without source ground truth, discovery volume is not the same as correctness
or exact-version authority.\\

Open-source or source-known validation &
Whether recovered outputs match dependency evidence available outside the APK &
Strong for audited cases, but it is a validation lane, not proof that the
same APK-only method is generally correct.\\

Cross-tool benchmark or assessment &
Which tools fail under shared criteria and where methodology is inconsistent &
Useful for gap analysis, but metrics remain bounded by the benchmark's
dataset and task definitions.\\

Security or vulnerability mapping &
Whether detected versions can be connected to vulnerability data &
The vulnerability conclusion is valid only as far as the component identity,
exact version, and vulnerability database mapping are valid.\\
\end{longtable}
\endgroup

This taxonomy explains why the chapter does not pool reported performance
numbers. LibScout's profile precision, ORLIS's obfuscation precision, LibD's
market-scale discovery volume, ATVHUNTER's version-level results,
LibAttention's R8-resilient presence detection, SAD's class-structural
version scores, and TPL-Benchmark's cross-tool results answer different
questions
\cite{backes2016libscout,li2017libd,wang2018orlis,zhan2021atvhunter,pan2025libattention,zhou2025sad,gu2025tplbenchmark}.
For a reconstructed \SBOM, the evaluation question must be restated as:
\emph{which claim type is being validated, against which denominator, and
under which transformations?}

\section{Comparative Analysis}

\Cref{tab:master-tier-a-comparison} applies the chapter taxonomy to the Tier
A seed corpus. The final column states the trust interpretation rather than a
simple score. This is deliberate: the same study can be strong evidence for
one claim type and weak evidence for another. The table records only supported
roles and limitations; missing study fields are not filled by inference.

\begingroup
\small
\renewcommand{\arraystretch}{1.15}
\begin{longtable}{p{0.18\textwidth}p{0.18\textwidth}p{0.26\textwidth}p{0.26\textwidth}}
\caption{Tier A studies classified by recovery role and trust implication.}
\label{tab:master-tier-a-comparison}\\
\toprule
\textbf{Study} & \textbf{Primary role} & \textbf{Main evidence and technique} & \textbf{SBOM-trust implication}\\
\midrule
\endfirsthead
\toprule
\textbf{Study} & \textbf{Primary role} & \textbf{Main evidence and technique} & \textbf{SBOM-trust implication}\\
\midrule
\endhead
\bottomrule
\endfoot
LibScout \cite{backes2016libscout} &
Reference-profile identity and version matching &
Fuzzy method descriptors, package-class-method profile hierarchy, similarity
score &
Supports the idea of version profiles, but package and method-structure
assumptions must be exposed as evidence limits.\\

LibD \cite{li2017libd} &
Scalable library discovery and boundary recovery &
Homogeny graph, call/inheritance/inclusion relations, opcode-derived
fingerprints &
Strengthens identity evidence under name obfuscation, but does not provide
exact-version or SBOM authority.\\

LibPecker \cite{zhang2018libpecker} &
Class-dependency based detection &
Class dependency hashing, adaptive similarity, weighted class matching &
Shows that structural class relations are useful, while similarity still
needs explicit thresholds and uncertainty handling.\\

ORLIS \cite{wang2018orlis} &
Obfuscation-resilient structural matching &
Call-graph-derived coarse features with \texttt{sdhash}, followed by
\texttt{ssdeep} class signatures and injective class mapping &
Provides partial-inclusion evidence, but high cost and missing exact-version
authority limit direct SBOM use.\\

LibID \cite{zhang2019libid} &
Obfuscated library and version identification &
Consistent class correspondence through binary integer programming, with
source-backed threshold evaluation &
Supports specific-version claims under its tested transformations; downstream
security authority still depends on reference and vulnerability evidence.\\

ASE 2020 assessment \cite{zhan2020assessment} &
Cross-tool failure analysis &
Unified criteria, tool patches, multidex and AAR observations &
Defines required evaluation dimensions and warns that engineering blind spots
can dominate measured accuracy.\\

Systematic assessment \cite{zhan2021systematic} &
Broader comparative methodology &
Six evaluation criteria, module-decoupling comparison, tool availability and
limitations &
Supports taxonomy-based synthesis and shows that CDG-style evidence is
preferable to package-only evidence.\\

ATVHUNTER \cite{zhan2021atvhunter} &
Version and vulnerability-oriented detection &
CDG decoupling, CFG coarse filter, opcode fuzzy matching, vulnerability DB &
Strongly connects version evidence to security use, but still lacks standard
SBOM provenance and full R8-inlining coverage.\\

LibScan \cite{wu2023libscan} &
Precise TPL identification &
Class correspondence, method opcode similarity, call-chain evidence &
Improves fine-grained matching evidence; claim strength remains tied to
reference coverage and transformation conditions.\\

Ramirez thesis \cite{ramirez2024empirical} &
Cross-platform tool comparison &
ORLIS and LIBLOOM Android datasets, manually selected Google Play apps, and
Android/iOS contrast &
Supports cross-tool and dataset-reuse analysis, but does not provide
source-known Android ground truth or standard SBOM validation.\\

LibHunter \cite{xie2024libhunter} &
Optimization-aware detection &
Call-site-optimization feature enhancement and cross-inlining method matching &
Makes optimization a first-class validity condition for exact-version claims.\\

LibAttention \cite{pan2025libattention} &
Learning-based presence detection &
Transformer over Smali tokens and SDK structural vocabulary &
Promising for obfuscation-resilient identity hints, but insufficient alone
for exact versions, provenance, or CVE authority.\\

SAD \cite{zhou2025sad} &
Class-structural version-level detection &
CDG neighbor aggregation, stateful/stateless split, field-operation sequence,
literal tie-breakers &
Provides strong version-discriminating evidence, while remaining sensitive to
method inlining and lacking standard SBOM output.\\

TPL-Benchmark \cite{gu2025tplbenchmark} &
Ground-truth benchmark and cross-tool reassessment &
Version-annotated real applications, dependency extraction, and standardized
evaluation of ten detectors &
Provides realistic comparative evidence and exposes low or uneven
generalization; it diagnoses detector limits rather than supplying a detector.\\

UniBOM \cite{safronov2025unibom} &
Adjacent SBOM generation and analysis workflow &
Firmware extraction, filesystem and source scanning, SBOM generation, CVE
analysis, and visualization &
Shows end-to-end SBOM integration outside Android TPL attribution, but cannot
justify APK Maven identity or exact-version claims.\\
\end{longtable}
\endgroup

Three patterns emerge from this comparison. First, the strongest exact-version
approaches are multi-evidence approaches: they do not rely on names alone but
combine structural, opcode, class, and literal information. Second,
evaluation authority is not the same as detection authority. Assessment
studies and source-known validation do not identify libraries by themselves,
but they determine whether a result can be interpreted safely. Third,
standard SBOM trust requirements are largely external to the detector output:
provenance, ambiguity, validation class, relationship type, and vulnerability
authority are rarely represented as first-class fields.

\section{Limitations of Existing Research}

The taxonomy exposes limitations that recur across the verified literature:

\begin{enumerate}
  \item \textbf{Library identity and exact version are often conflated.}
  Presence detection, class-set matching, and exact-version attribution are
  different claims, but tool outputs do not always preserve the distinction.

  \item \textbf{Reference coverage is rarely treated as part of the result.}
  A missed match can result from absent references, extraction failure, or a
  weak matcher. Without coverage metadata, the user cannot tell which failure
  occurred.

  \item \textbf{Transformation coverage is incomplete.}
  Identifier renaming is widely considered, but shrinking, partial inclusion,
  package flattening, and method inlining remain unevenly evaluated
  \cite{xie2024libhunter,pan2025libattention,zhou2025sad}.

  \item \textbf{Confidence is usually a threshold outcome, not an auditable
  explanation.}
  Many methods compute similarities or classifier scores, but few expose why
  a claim should be trusted, which signals conflicted, or why the method did
  not abstain.

  \item \textbf{Partial libraries and dependency relationships are not fully
  represented.}
  Similarity scores can detect partial inclusion, but a standard SBOM also
  needs relationship semantics: direct, transitive, bundled, shaded, copied,
  or uncertain.

  \item \textbf{Vulnerability interpretation is stronger than the underlying
  evidence in many pipelines.}
  A detected library name is not enough for a version-specific \CVE claim.
  Even a plausible version candidate requires reference, transformation, and
  validation context.

  \item \textbf{Reproducibility remains uneven.}
  Tool availability, installation difficulty, dataset access, and benchmark
  construction vary across studies, limiting independent confirmation of
  reported results.

  \item \textbf{Standard SBOM output is not the dominant research artifact.}
  The verified studies mostly produce detector outputs rather than complete
  CycloneDX or SPDX statements with evidence, provenance, and uncertainty
  fields.
\end{enumerate}

These limitations do not make the existing work weak. Rather, they show that
the literature has mostly optimized detection performance, while trustworthy
\SBOM recovery requires an additional layer that explains what each recovered
claim means and how it may be used.

\section{Research Gap}

The state of the art provides valuable techniques for recognizing Android
third-party libraries from compiled artifacts. Reference profiles, fuzzy
descriptors, dependency graphs, opcode features, call-chain evidence,
stateful/stateless class summaries, learned Smali representations, and
source-known validation all contribute useful evidence. Within the reviewed
Android TPL literature and adjacent binary-SBOM studies, however, no verified
study provides a unified model that connects these signals to a reconstructed
\SBOM claim while preserving all of the trust dimensions used in this thesis.

The research gap can therefore be stated as follows:

\begin{quote}
Existing Android component detectors provide useful identification and
version-discrimination techniques, but they do not consistently attach each
recovered SBOM claim to its evidence provenance, reference coverage,
transformation risk, confidence boundary, ambiguity, validation class, and
permitted vulnerability interpretation.
\end{quote}

This gap is not only a formatting problem. A syntactically valid SBOM can
still be untrustworthy if it hides uncertainty, promotes identity evidence
into exact-version authority, or turns an ambiguous version set into a
definitive \CVE finding. Conversely, a detector that abstains or reports a
bounded version set may be more useful for security decision-making than a
pipeline that always emits a single version without explaining the evidence.

\section{Conclusion}

This chapter classified Android component-recovery research by technique,
evidence source, transformation threat, and evaluation practice. The synthesis
shows a clear technical progression from package and profile matching toward
graph-based, version-aware, optimization-aware, and learning-based methods.
It also shows a persistent trust gap: detector outputs rarely preserve enough
provenance, uncertainty, validation context, and security authority to become
auditable reconstructed \SBOM claims on their own.

The general conclusion returns to this gap through the research questions and
future perspectives. The key lesson is that the reviewed techniques should not
be judged only by whether they emit a name or version. They should also be
judged by whether their outputs communicate why the recovered component entry
should be believed, where its evidence stops, and how safely it may be used.
